This project connects standard Enterprise Architecture modeling with Semantic Web-based analysis. By formalizing a translation pipeline from ArchiMate XML to an RDF/OWL Knowledge Graph, we established a semantic basis for architectural analysis. The implementation shows how knowledge enrichment---such as automating TOGAF Gap Analysis and identifying EA smells---can be executed directly within the triplestore via SPARQL updates to support the tracking of architectural evolution.

Furthermore, our LLM evaluation yields preliminary insights into natural-language model exploration. In our setup, the interactive GraphDB condition scored higher than the static XML, JSON, and Turtle conditions. This comparison is limited by the two runs per condition and the incomplete 2x2 design, so it should not be interpreted as evidence that access mode generally supersedes data format. The interactive approach also revealed challenges in grounding, as the generation model frequently described its query strategy rather than explicitly tracing the instance-level identifiers responsible for its conclusions.

While the results are encouraging, certain limitations must be acknowledged. First, the answers were initially scored using a dedicated LLM-based grader. To mitigate potential grading errors, we performed targeted manual verification on outliers and complex edge cases, making minor adjustments where necessary. However, a complete manual review of all generated answers was not feasible, meaning the evaluation still relies heavily on an automated metric. Second, the evaluation was limited to two runs per condition, meaning the observed performance gaps between serialization formats lack statistical significance. Finally, the semantic queries were tested primarily against the standardized ArchiSurance model.

Future work should address these limitations by scaling the experimental runs to establish robust statistical baselines and validating the pipeline against large-scale, real-world enterprise repositories. Technical extensions could also incorporate structural graph analytics natively within the graph to further aid architectural decision-making.
